\begin{table}[H]
    \centering
    \linespread{1.05}\selectfont
    \caption{Moduły programowe symulatora wraz z~ich odpowiedzialnością, interfejsem i~odniesieniem do teorii}
    \label{tab:software_modules}
    \scriptsize
    \begin{tabular}{|p{3.3cm}|p{4.5cm}|p{4.3cm}|p{2.2cm}|}
        \hline
        \textbf{Plik} & \textbf{Główna odpowiedzialność} & \textbf{Kluczowe wejścia $\rightarrow$ wyjścia} & \textbf{Teoria} \\
        \hline
        \multicolumn{4}{|l|}{\textit{Fasady potoku (\texttt{src/})}} \\
        \hline
        \texttt{step1\_\allowbreak transmitter.m} & Synteza ramki PPDU IEEE~802.11a/ax z~losowym polem danych & \texttt{WIFI\_STANDARD}, \texttt{cbw}, \texttt{mcs}, \texttt{psdu\_length} $\rightarrow$ \texttt{waveform}, \texttt{fs}, \texttt{params}, \texttt{cfg} & \ref{subsec:wifi_signal}, tab.~\ref{tab:wifi_params} \\
        \hline
        \texttt{step2\_\allowbreak channel.m} & Kanał TGax (DPI i~echa stałe), synteza ech celów ruchomych, szum termiczny wg budżetu mocy & \texttt{waveform}, \texttt{fs}, \texttt{model}, \texttt{cfg}, \texttt{targets}, \texttt{p\_rx\_dBm}, \texttt{nf\_dB} $\rightarrow$ \texttt{signal}, \texttt{fc}, \texttt{intrinsic\_delay} & \ref{subsec:pbr_fundamentals} \\
        \hline
        \texttt{step3\_\allowbreak receiver.m} & Synchronizacja ramki, demodulacja, estymacja i~kondycjonowanie transmitancji, mapa odległość--prędkość & \texttt{signal}, \texttt{waveform}, \texttt{fs}, \texttt{fc}, \texttt{params}, \texttt{sync\_idx} $\rightarrow$ \texttt{radar\_data} & \ref{subsec:ofdm_modulation}, \ref{subsec:channel_estimation}--\ref{subsec:range_doppler_map} \\
        \hline
        \texttt{step4\_\allowbreak interpreter.m} & Wybór detektora, progowanie CFAR, ekstrakcja celów, metryka PNFR & \texttt{radar\_data}, \texttt{p\_FA}, \texttt{stc\_method} $\rightarrow$ \texttt{targets}, \texttt{CPer\_F}, \texttt{Per\_F}, \texttt{eta}, \texttt{pnfr}, \texttt{diag\_info} & \ref{subsec:target_detection_extraction} \\
        \hline
        \multicolumn{4}{|l|}{\textit{Parametryzacja nadajnika i~kanału (\texttt{lib/})}} \\
        \hline
        \texttt{get\_wifi\_\allowbreak params.m} & Obiekt \texttt{cfg}, $N_{\mathrm{FFT}}$, $N_{\mathrm{CP}}$, automatyczne skalowanie długości ładunku, maska aktywnych podnośnych, indeksy luk DC & \texttt{standard}, \texttt{psdu\_length}, \texttt{cbw}, \texttt{mcs} $\rightarrow$ \texttt{params} & \ref{subsec:wifi_signal}, tab.~\ref{tab:wifi_params} \\
        \hline
        \texttt{get\_carrier\_\allowbreak frequency.m} & Odwzorowanie nazwy pasma na częstotliwość nośną & \texttt{band\_input} $\rightarrow$ \texttt{fc}, \texttt{band\_str} & \ref{subsec:pbr_fundamentals} \\
        \hline
        \texttt{configure\_\allowbreak targets.m} & Przeliczenie opisu fizycznego celu na opóźnienie, przesunięcie Dopplera i~amplitudę echa & \texttt{targets}, \texttt{fs}, \texttt{fc}, \texttt{channel\_model}, \texttt{p\_rx\_dBm} $\rightarrow$ \texttt{delays}, \texttt{dopplers}, \texttt{taps} & \ref{subsec:pbr_fundamentals} \\
        \hline
        \texttt{get\_tgax\_\allowbreak propagation\_\allowbreak params.m} & Parametry dwuodcinkowego modelu tłumienia TGn/TGax (odległość przełamania, wykładniki LOS/NLOS) & \texttt{model} $\rightarrow$ \texttt{p} (\texttt{d\_bp}, \texttt{n\_los}, \texttt{n\_nlos}) & \ref{subsec:channel_implementation} \\
        \hline
        \multicolumn{4}{|l|}{\textit{Przetwarzanie sygnału w~odbiorniku (\texttt{lib/})}} \\
        \hline
        \texttt{demodulate.m} & Usunięcie prefiksu cyklicznego i~kolumnowa transformata FFT & \texttt{signal}, \texttt{n\_fft}, \texttt{n\_cp} $\rightarrow$ \texttt{F\_matrix}, \texttt{n\_syms} & \ref{subsec:ofdm_modulation} \\
        \hline
        \texttt{estimate\_\allowbreak channel.m} & Estymacja ZF, współczynnik $\zeta_{\mathrm{ZF}}$, filtr MTI w~czasie wolnym, wywołanie interpolacji luki DC & \texttt{F\_rx}, \texttt{F\_tx}, \texttt{params} $\rightarrow$ \texttt{H\_shifted}, \texttt{H\_raw}, \texttt{zeta\_dB}, \texttt{H\_no\_static} & \ref{subsec:channel_estimation}, \ref{subsec:clutter_filtering} \\
        \hline
        \texttt{interpolate\_\allowbreak dc\_nulls.m} & Liniowa interpolacja spójnych luk zerowych podnośnych wewnętrznych & \texttt{H}, \texttt{dc\_indices} $\rightarrow$ \texttt{H} & \ref{subsec:clutter_filtering} \\
        \hline
        \texttt{compute\_range\_\allowbreak doppler\_map.m} & Okno 2D Blackmana--Harrisa, dopełnianie zerami, periodogram zespolony i~periodogram mocy & \texttt{H\_shifted}, \texttt{Nfft}, \texttt{n\_symbols}, \texttt{active\_mask} $\rightarrow$ \texttt{CPer\_F}, \texttt{W}, \texttt{Per\_F}, \texttt{N\_per}, \texttt{M\_per} & \ref{subsec:range_doppler_map} \\
        \hline
        \texttt{calculate\_\allowbreak physical\_axes.m} & Odwzorowanie indeksów DFT/IDFT na osie odległości i~prędkości & \texttt{fs}, \texttt{fc}, \texttt{Nfft}, \texttt{Ncp}, \texttt{N\_per}, \texttt{M\_per} $\rightarrow$ \texttt{axis\_range}, \texttt{axis\_velocity} & \ref{subsec:pbr_fundamentals}, \ref{subsec:range_doppler_map} \\
        \hline
        \multicolumn{4}{|l|}{\textit{Detekcja i~ocena wyników (\texttt{lib/})}} \\
        \hline
        \texttt{get\_detection\_\allowbreak algorithm.m} & Centralny rejestr detektorów (identyfikator, nazwa, uchwyt wykonawczy) & \texttt{method} $\rightarrow$ \texttt{alg} & \ref{subsec:target_detection_extraction} \\
        \hline
        \texttt{binary\_\allowbreak successive\_\allowbreak target\_\allowbreak cancellation.m} & Próg CFAR, iteracyjna detekcja z~maskowaniem binarnym listka głównego & \texttt{Per\_F}, \texttt{W}, \texttt{p\_FA} $\rightarrow$ \texttt{targets}, \texttt{B}, \texttt{eta}, \texttt{diag\_info} & \ref{subsec:target_detection_extraction} \\
        \hline
        \texttt{coherent\_\allowbreak successive\_\allowbreak target\_\allowbreak cancellation.m} & Próg CFAR, iteracyjna detekcja z~koherentnym odejmowaniem repliki celu & \texttt{CPer\_F}, \texttt{Per\_F}, \texttt{W}, \texttt{p\_FA} $\rightarrow$ \texttt{targets}, \texttt{CPer\_F}, \texttt{eta}, \texttt{diag\_info} & \ref{subsec:target_detection_extraction} \\
        \hline
        \texttt{compute\_pnfr.m} & Kontrast szczytu celu względem tła szumowo-zakłóceniowego (PNFR) & \texttt{CPer}, \texttt{target\_idx}, \texttt{guard} $\rightarrow$ \texttt{metrics} & \ref{subsec:target_detection_extraction} \\
        \hline
    \end{tabular}
\end{table}
